% Created 2025-02-17 lun 14:14
% Intended LaTeX compiler: pdflatex
\documentclass[11pt]{article}
\usepackage[utf8]{inputenc}
\usepackage[T1]{fontenc}
\usepackage{graphicx}
\usepackage{longtable}
\usepackage{wrapfig}
\usepackage{rotating}
\usepackage[normalem]{ulem}
\usepackage{amsmath}
\usepackage{amssymb}
\usepackage{capt-of}
\usepackage{hyperref}
\usepackage{minted}

\renewcommand{\contentsname}{Tabla de contenidos}
\usepackage{parskip}
\usepackage[utf8]{inputenc} %% For unicode chars
\usepackage{placeins}
\usepackage[margin=2.50cm]{geometry}
\usepackage[T1]{fontenc}
\usepackage{mathpazo}
\linespread{1.05}
\usepackage[scaled]{helvet}
\usepackage{courier}
\hypersetup{colorlinks=true,linkcolor=blue}
\RequirePackage{fancyvrb}
\DefineVerbatimEnvironment{verbatim}{Verbatim}{fontsize=\small,formatcom = {\color[rgb]{0.5,0,0}}}
\usepackage{mdframed}
\BeforeBeginEnvironment{minted}{\begin{mdframed}}
\AfterEndEnvironment{minted}{\end{mdframed}}
\setcounter{secnumdepth}{2}
\date{}
\title{Despliegue api rest laravel 10}
\hypersetup{
 pdfauthor={},
 pdftitle={Despliegue api rest laravel 10},
 pdfkeywords={},
 pdfsubject={},
 pdfcreator={Emacs 27.1 (Org mode 9.6.18)}, 
 pdflang={Spanish}}
\begin{document}

\maketitle
\setcounter{tocdepth}{2}
\tableofcontents

\setlength\parindent{10pt}

\section{Despliegue de API REST en Laravel 10 con Sail en Producción}
\label{sec:orgf9dd177}

Vamos a desplegar tu API REST en Laravel 10 con Docker y Sail en un
servidor de producción. Te explicaré cada paso detalladamente.

\subsection{Configuración del Servidor}
\label{sec:org3932421}

Necesitas un VPS o servidor dedicado con un sistema operativo Linux
(ej. Ubuntu 22.04). Lo primero es actualizarlo y preparar el entorno.

\subsubsection*{En el Servidor (host)}
\label{sec:orgb0f901b}
Conéctate vía SSH:

\begin{minted}[]{sh}
ssh usuario@tu-servidor
\end{minted}

Actualiza el sistema:

\begin{minted}[]{sh}
sudo apt update && sudo apt upgrade -y
\end{minted}

Instala los paquetes básicos:

\begin{minted}[]{sh}
sudo apt install -y curl git unzip ufw
\end{minted}

\subsection{Instalar Docker y Docker Compose}
\label{sec:orge90b4ff}

Sail funciona sobre Docker, así que necesitas instalarlo en el
servidor.

\subsubsection*{En el Servidor (host)}
\label{sec:orgfc854c2}
Ejecuta:

\begin{minted}[]{sh}
sudo apt install -y docker.io docker-compose
sudo systemctl enable --now docker
\end{minted}

Verifica que Docker funciona:

\begin{minted}[]{sh}
docker --version
docker-compose --version
\end{minted}

Añade tu usuario al grupo Docker (opcional, para no usar `sudo` en
cada comando):

\begin{minted}[]{sh}
sudo usermod -aG docker $(whoami)
\end{minted}

\textbf{\textbf{(Sal de la sesión y vuelve a entrar para que surta efecto)}}

\subsection{Subir el Código de la API}
\label{sec:orgc9bc209}

Tienes dos opciones: subir el código desde tu máquina o clonar el
repositorio en el servidor.

\subsubsection*{Opción 1: Subir desde tu máquina}
\label{sec:org1543641}
En tu máquina local:

\begin{minted}[]{sh}
scp -r tu-api usuario@tu-servidor:/var/www/
\end{minted}

\subsubsection*{Opción 2: Clonar directamente en el servidor}
\label{sec:org4911543}
En el servidor:

\begin{minted}[]{sh}
cd /var/www/
git clone https://github.com/tu-usuario/tu-api.git
\end{minted}

Asegúrate de que los permisos sean correctos:

\begin{minted}[]{sh}
sudo chown -R $USER:www-data /var/www/tu-api
sudo chmod -R 775 /var/www/tu-api
\end{minted}

\subsection{Configurar Laravel y Entorno}
\label{sec:org3145d40}

\subsubsection*{Dentro del Proyecto en el Servidor}
\label{sec:orge6c8b9d}
Accede a la carpeta del proyecto:

\begin{minted}[]{sh}
cd /var/www/tu-api
\end{minted}

Copia el archivo de entorno y edítalo:

\begin{minted}[]{sh}
cp .env.example .env
nano .env
\end{minted}

Cambia las siguientes variables:

\begin{minted}[]{bash}
APP_ENV=production
APP_DEBUG=false
APP_URL=https://tu-dominio.com

DB_CONNECTION=mysql
DB_HOST=mysql
DB_PORT=3306
DB_DATABASE=nombre_bd
DB_USERNAME=usuario_bd
DB_PASSWORD=clave_secreta
\end{minted}

\subsection{Configurar Sail para Producción}
\label{sec:orge8ef64e}

\textbf{\textbf{(En `/var/www/tu-api` en el servidor)}}

Levanta los contenedores:

\begin{minted}[]{sh}
./vendor/bin/sail up -d
\end{minted}

Ejecuta las migraciones:

\begin{minted}[]{sh}
./vendor/bin/sail artisan migrate --force
\end{minted}

Genera la clave de la aplicación:

\begin{minted}[]{sh}
./vendor/bin/sail artisan key:generate
\end{minted}

\subsection{Optimización de Laravel}
\label{sec:orgdbbde4b}

Dentro del contenedor de Laravel:

\begin{minted}[]{sh}
./vendor/bin/sail artisan config:cache
./vendor/bin/sail artisan route:cache
./vendor/bin/sail artisan view:cache
\end{minted}

\subsection{Configurar Nginx como Proxy Reverso}
\label{sec:org3d9b2a3}

Sail ya tiene Nginx dentro del contenedor, pero necesitas un \textbf{\textbf{proxy
inverso}} para que la API sea accesible desde el exterior.

\subsubsection*{En el Servidor (host)}
\label{sec:orgefb721d}
Instala Nginx:

\begin{minted}[]{sh}
sudo apt install -y nginx
\end{minted}

Crea una configuración para Laravel:

\begin{minted}[]{sh}
sudo nano /etc/nginx/sites-available/tu-api
\end{minted}

Pega este contenido:

\begin{minted}[]{nginx}
server {
    listen 80;
    server_name tu-dominio.com;

    location / {
        proxy_pass http://127.0.0.1:80;
        proxy_set_header Host $host;
        proxy_set_header X-Real-IP $remote_addr;
        proxy_set_header X-Forwarded-For $proxy_add_x_forwarded_for;
        proxy_set_header X-Forwarded-Proto $scheme;
    }
}
\end{minted}

Activa el sitio y reinicia Nginx:

\begin{minted}[]{sh}
sudo ln -s /etc/nginx/sites-available/tu-api /etc/nginx/sites-enabled/
sudo systemctl restart nginx
\end{minted}

\subsection{Configurar HTTPS con Let's Encrypt}
\label{sec:orge04e614}

Ejecuta:

\begin{minted}[]{sh}
sudo apt install certbot python3-certbot-nginx
sudo certbot --nginx -d tu-dominio.com
\end{minted}

Esto generará y configurará automáticamente el certificado SSL.

\subsection{Configurar Supervisor para las Colas (Opcional)}
\label{sec:org6c48a73}

Si tu aplicación usa colas, instala `supervisor` y configúralo.

\subsubsection*{En el Servidor (host)}
\label{sec:orgaf8c1d4}
\begin{minted}[]{sh}
sudo apt install -y supervisor
\end{minted}

Crea el archivo:

\begin{minted}[]{sh}
sudo nano /etc/supervisor/conf.d/laravel-worker.conf
\end{minted}

Agrega esto:

\begin{minted}[]{ini}
[program:laravel-worker]
process_name=%(program_name)s_%(process_num)02d
command=docker exec -i tu-api-app sail artisan queue:work --tries=3
autostart=true
autorestart=true
numprocs=1
redirect_stderr=true
stdout_logfile=/var/log/laravel-worker.log
\end{minted}

Carga la configuración:

\begin{minted}[]{sh}
sudo supervisorctl reread
sudo supervisorctl update
sudo supervisorctl start laravel-worker
\end{minted}

\subsection{Seguridad y Firewall}
\label{sec:org1821faa}

Habilita el firewall y solo permite SSH, HTTP y HTTPS:

\begin{minted}[]{sh}
sudo ufw allow OpenSSH
sudo ufw allow 80
sudo ufw allow 443
sudo ufw enable
\end{minted}

\subsection{Últimos Pasos y Pruebas}
\label{sec:org686f978}

\subsubsection*{🔹 En el Servidor (host)}
\label{sec:orga8cb73c}
Verifica que todo está corriendo:

\begin{minted}[]{sh}
docker ps
\end{minted}

Si ves los contenedores en ejecución, ¡vas bien!

Prueba que la API responde:

\begin{minted}[]{sh}
curl -I https://tu-dominio.com/api/endpoint
\end{minted}

Si necesitas ver logs:

\begin{minted}[]{sh}
./vendor/bin/sail artisan logs
\end{minted}

\subsection{Resumen Final}
\label{sec:org05a8fa0}
\begin{itemize}
\item \textbf{\textbf{Servidor}}: Instalamos Docker, Docker Compose y Nginx.
\item \textbf{\textbf{Proyecto}}: Subimos Laravel y configuramos `.env`.
\item \textbf{\textbf{Sail}}: Levantamos los contenedores (`sail up -d`).
\item \textbf{\textbf{Nginx}}: Configuramos un proxy inverso.
\item \textbf{\textbf{SSL}}: Instalamos Let's Encrypt.
\item \textbf{\textbf{Supervisor}}: Configuramos workers si usamos colas.
\item \textbf{\textbf{Seguridad}}: Configuramos firewall y permisos.
\end{itemize}

\textbf{\textbf{¡Tu API REST en Laravel 10 con Sail ya está en producción!}}
\end{document}
